\chapter{Machine-Learning Algorithms}
\label{ch36:ml}

In Chapter~35 we learned how to distribute computation across
many cores at once, parallelizing loops, matrix multiplications, and
divide-and-conquer combines. Throughout those chapters, we wrote algorithms that
\emph{know} what to do: you told sorting how to compare, you told Dijkstra how
edges relax, and you told dynamic programming exactly which subproblems to
solve. Every line of pseudocode was a rule you supplied.

This chapter takes a different approach. Here we will write algorithms that
do \emph{not} start out knowing what to do. Instead, we hand them a pile of
examples and a way to measure how badly they are doing, and we let them
\emph{adjust themselves} until they do well. The rule is no longer something
you write down by hand; it is something the algorithm \emph{finds}.

The marketing around ``machine learning'' can make this sound mysterious, but
the three core ideas of this chapter -- clustering, gradient descent, and the
perceptron -- are ordinary algorithms in the exact sense of everything you
have already read. They have inputs, outputs, loops, invariants, and running
times. The only genuinely new idea is that the \emph{objective} -- the thing we
are optimizing -- is now a quantity defined by data, and ``solving the problem''
means ``making that quantity small.''

\begin{focusbox}
The main idea of the chapter:
\emph{learning is optimization over data.} You define a number that measures
``how wrong am I'' (a \textbf{loss} or \textbf{cost}), and then you run an
algorithm that pushes that number down. Everything else -- k-means, gradient
descent, the perceptron -- is a specific choice of \emph{what} number and
\emph{how} to push it down.
\end{focusbox}

\begin{historybox}
The idea of machines that learn from examples rather than fixed instructions
goes back to the early years of postwar computing. In 1957, Frank Rosenblatt at
the Cornell Aeronautical Laboratory built the \emph{Mark~I Perceptron}, an early
attempt to instantiate biological neural computation in custom electronic hardware.
The machine was an imposing room-sized apparatus: a camera with a $20\times 20$ grid
of cadmium-sulfide photocells served as the retina, connected to hundreds of
motor-driven potentiometers (variable resistors that adjust electrical current)
that physically rotated to adjust their analog weights during training. Rosenblatt
proved the Perceptron Convergence Theorem (formalized by Albert Novikoff in 1962),
establishing that if a linearly separable boundary exists, the machine is
guaranteed to find it in a finite number of mistake updates.

That same year, in 1957, Stuart Lloyd at Bell Telephone Laboratories was grappling
with pulse-code modulation (PCM)---converting continuous analog human voice
waveforms into discrete digital pulses for transmission over telephone wires.
To minimize distortion, Lloyd needed to quantize continuous signal voltages into a
finite set of representative voltage levels. He proposed an alternating algorithm:
partition the voltage range into nearest-neighbor cells, then update each code
level to the centroid (the average value) of its cell. That memorandum became the
classic $k$-means algorithm (formally published in 1982), which remains the standard
method for unsupervised clustering across science and industry.

Closely related to both the perceptron and $k$-means is an optimization method
first described by Augustin-Louis Cauchy in 1847: the method of steepest
descent, taking small steps opposite the gradient to minimize a function. Yet the
field was nearly extinguished in 1969 when Marvin Minsky and Seymour Papert
published \emph{Perceptrons}, mathematically proving that a single-layer perceptron
cannot even learn the simple non-linear \textsc{XOR} (exclusive OR) function. The
resulting disillusionment triggered the first ``AI Winter,'' freezing neural
network research funding for nearly two decades until David Rumelhart, Geoffrey
Hinton, and Ronald Williams demonstrated in 1986 that multi-layer networks trained
by backpropagation (calculus's chain rule applied to layers of weights) easily
learn non-linear representations. Today, learning algorithms drive automated speech,
vision, and scientific modeling, though open theoretical questions remain:
why heavily overparameterized deep networks generalize so well instead of
merely memorizing their training data, and how to guarantee safety on non-convex
loss surfaces riddled with saddle points.
\end{historybox}

\section{Learning as optimization: the one idea}
\label{ch36:learning-as-optimization}

Let us make the slogan concrete before we touch any specific method.

Imagine you are given a bunch of data points and a \emph{model}: a function
$f_\theta$ with some adjustable knobs collected into a vector $\theta$ (the
Greek letter ``theta,'' our standard name for ``the parameters''). Turning the
knobs changes the function. For each setting of $\theta$ you can ask: how well
does $f_\theta$ fit the data? You answer that with a single number, the
\textbf{cost} (or \textbf{loss}) $J(\theta)$. Low cost means good fit; high cost
means bad fit.

\begin{intuitionbox}
Picture $J$ as a curved surface. The horizontal directions are the knob settings
$\theta$; the height above each point is the cost $J(\theta)$. A \emph{good}
model is a \emph{low} point on this surface -- a valley. Learning is just
\emph{walking downhill}. You start somewhere, look around for a downward
direction, take a step, and repeat until the ground is flat under your feet.
That flat spot is your trained model.
\end{intuitionbox}

\begin{formalbox}
A \textbf{supervised learning problem} is specified by:
\begin{itemize}
  \item a dataset of $n$ examples $(x^{(1)},y^{(1)}),\dots,(x^{(n)},y^{(n)})$,
        where each $x^{(i)}\in\R^d$ is an \emph{input} (a vector of $d$
        \emph{features}) and each $y^{(i)}$ is a \emph{label};
  \item a model family $\{f_\theta : \theta\in\R^p\}$, indexed by a parameter
        vector $\theta$;
  \item a per-example loss $\ell(f_\theta(x), y)\ge 0$.
\end{itemize}
The training objective is the average loss over the data,
\[
  J(\theta) \;=\; \frac{1}{n}\sum_{i=1}^{n} \ell\!\big(f_\theta(x^{(i)}),\, y^{(i)}\big),
\]
and \emph{to train} means to (approximately) solve
$\theta^\star = \arg\min_{\theta} J(\theta)$.
In \emph{unsupervised} problems (like clustering) there are no labels $y$; the
cost is defined purely in terms of the inputs.
\end{formalbox}

That is the whole framework, visualized as a closed optimization loop
in Figure~\ref{ch36:fig-learning-loop}. Now we just need (1) ways to write down
useful costs, and (2) algorithms to minimize them. We will meet both in three
flavors.

\begin{figure}[h]
\centering
\begin{tikzpicture}[>=Stealth, font=\small]
  % Stage 1: Data
  \node[draw=gray!70, fill=gray!8, rounded corners=3pt, minimum width=2.2cm, minimum height=1.1cm, align=center] (data) at (0, 1.5)
    {\textbf{Training Data}\\$(x^{(i)}, y^{(i)})$};

  % Stage 2: Model
  \node[draw=sourceblue, fill=blue!8, rounded corners=3pt, minimum width=2.4cm, minimum height=1.1cm, align=center] (model) at (4.0, 1.5)
    {\textbf{Model } $f_\theta$\\Parameters $\theta$};

  % Stage 3: Loss & Cost
  \node[draw=quotientorange, fill=orange!8, rounded corners=3pt, minimum width=2.6cm, minimum height=1.1cm, align=center] (loss) at (8.8, 1.5)
    {\textbf{Loss \& Cost}\\Total error $J(\theta)$};

  % Stage 4: Optimizer
  \node[draw=targetgreen, fill=green!8, rounded corners=3pt, minimum width=2.8cm, minimum height=1.1cm, align=center] (optim) at (6.6, -1.2)
    {\textbf{Optimizer}\\$\theta \gets \theta - \eta\,\nabla J$};

  % Forward arrows
  \draw[->, thick] (data.east) -- (model.west)
    node[midway, above, font=\scriptsize] {inputs $x^{(i)}$};
  \draw[->, thick] (model.east) -- (loss.west)
    node[midway, above, font=\scriptsize] {predictions $\hat{y}$};

  % Ground truth label over the top (no crossing)
  \draw[->, thick, dashed, gray!70] (data.north) -- ++(0, 0.5) -| (loss.north)
    node[pos=0.25, above, font=\scriptsize, gray!80] {ground truth labels $y^{(i)}$};

  % Downward to optimizer
  \draw[->, thick] (loss.south) |- (optim.east)
    node[pos=0.5, right, font=\scriptsize] {gradient $\nabla J$};

  % Feedback to model
  \draw[->, very thick, targetgreen!80!black] (optim.west) -| (model.south)
    node[midway, left, font=\scriptsize, targetgreen!80!black] {update $\Delta\theta$};
\end{tikzpicture}
\caption{The supervised machine learning loop. Inputs $x^{(i)}$ flow forward through
the parameterized model $f_\theta$ to produce predictions $\hat{y}^{(i)}$. The loss
evaluates error against true labels $y^{(i)}$ and aggregates into total cost $J(\theta)$.
An optimizer computes gradients $\nabla J(\theta)$ and feeds updated parameters back
into the model, driving iterative descent.}
\label{ch36:fig-learning-loop}
\end{figure} This connects to earlier chapters: greedy algorithms
made a locally best choice each step; gradient descent makes a locally
downhill choice each step. Dynamic programming reused subproblem answers; here
we reuse the \emph{current} parameters as the starting point for the next
improvement. Optimization has appeared throughout the book.

\section{Clustering with $k$-means}
\label{ch36:kmeans}

We begin with an \emph{unsupervised} problem because it shows the
``define-a-cost-then-minimize-it'' pattern at its cleanest.

You are handed $n$ points in space and no labels at all -- just the points.
You suspect they fall into natural groups. Maybe they are customers and the
groups are buying habits; maybe they are pixels and the groups are colors.
You want to discover $k$ groups (\textbf{clusters}) and a representative point
for each (a \textbf{centroid}).

\begin{intuitionbox}
$k$-means answers two linked questions:
\begin{enumerate}
  \item \emph{If I knew the centroids}, which cluster does each point belong to?
        Easy: assign each point to its \emph{nearest} centroid.
  \item \emph{If I knew the assignments}, where should each centroid sit?
        Easy: at the \emph{average} (the ``mean'') of its assigned points.
\end{enumerate}
Neither is known up front, so we guess, then alternate: re-assign, re-center,
re-assign, re-center. Each round can only make things tidier. That alternation
is the whole algorithm.
\end{intuitionbox}

Figure~\ref{ch36:fig-kmeans-loop} traces this alternating optimization on three
natural clusters, showing the centroids shifting to cluster means and the Voronoi
cell boundaries settling into place.

\begin{figure}[H]
\centering
\begin{tikzpicture}[
  scale=0.78,
  >=Stealth,
  ptA/.style={circle, fill=sourceblue!70, inner sep=1.2pt},
  ptB/.style={circle, fill=targetgreen!70, inner sep=1.2pt},
  ptC/.style={circle, fill=quotientorange!80, inner sep=1.2pt},
  cent/.style={circle, draw=black, very thick, inner sep=2pt, font=\scriptsize\bfseries},
  lbl/.style={font=\sffamily\scriptsize}
]
% Panel (a): Step 1 - Initial guess and Voronoi cells
\begin{scope}[shift={(0,0)}]
  % Cluster A points around (1, 3)
  \foreach \p in {(0.8,3.2),(1.2,2.8),(0.6,2.7),(1.4,3.3),(1.0,3.5)} \node[ptA] at \p {};
  % Cluster B points around (3.5, 3.2)
  \foreach \p in {(3.2,3.4),(3.8,3.1),(3.3,2.9),(3.7,3.5),(3.5,2.7)} \node[ptB] at \p {};
  % Cluster C points around (2.2, 0.9)
  \foreach \p in {(2.0,1.1),(2.5,0.8),(1.8,0.7),(2.3,1.3),(2.1,0.5)} \node[ptC] at \p {};

  % Voronoi partition lines (rough)
  \draw[dashed, gray!70] (0, 1.8) -- (4.2, 1.8);
  \draw[dashed, gray!70] (2.3, 1.8) -- (2.3, 4.0);

  % Initial random centroids
  \node[cent, fill=sourceblue!25] (m1) at (1.5, 2.2) {$\mu_1$};
  \node[cent, fill=targetgreen!25] (m2) at (3.2, 2.1) {$\mu_2$};
  \node[cent, fill=quotientorange!25] (m3) at (2.0, 2.8) {$\mu_3$};
  \draw[dashed, gray!70] (2.3, 2.0) -- (2.3, 4.0);

  \node[lbl, align=center] at (2.1, -0.6) {\textbf{(a) Initial centroids}\\Arbitrary guesses $\mu_j$\\Voronoi cells assign points};
\end{scope}

% Panel (b): Step 2 - Centroid shift
\begin{scope}[shift={(5.6,0)}]
  % Points
  \foreach \p in {(0.8,3.2),(1.2,2.8),(0.6,2.7),(1.4,3.3),(1.0,3.5)} \node[ptA] at \p {};
  \foreach \p in {(3.2,3.4),(3.8,3.1),(3.3,2.9),(3.7,3.5),(3.5,2.7)} \node[ptB] at \p {};
  \foreach \p in {(2.0,1.1),(2.5,0.8),(1.8,0.7),(2.3,1.3),(2.1,0.5)} \node[ptC] at \p {};

  % Old centroids (faint)
  \node[circle, draw=gray!40, inner sep=1.8pt] (old1) at (1.5, 2.2) {};
  \node[circle, draw=gray!40, inner sep=1.8pt] (old2) at (3.2, 2.0) {};
  \node[circle, draw=gray!40, inner sep=1.8pt] (old3) at (2.0, 2.8) {};

  % New centroids (cluster means)
  \node[cent, fill=sourceblue!40] (new1) at (1.0, 3.1) {$\mu_1$};
  \node[cent, fill=targetgreen!40] (new2) at (3.5, 3.1) {$\mu_2$};
  \node[cent, fill=quotientorange!40] (new3) at (2.1, 0.9) {$\mu_3$};

  % Shift vectors
  \draw[->, very thick, trapred!80!black] (old1) -- (new1);
  \draw[->, very thick, trapred!80!black] (old2) -- (new2);
  \draw[->, very thick, trapred!80!black] (old3) -- (new3);

  \node[lbl, align=center] at (2.1, -0.6) {\textbf{(b) Update centroids}\\Move each $\mu_j$ to mean\\of its assigned cluster};
\end{scope}

% Panel (c): Step 3 - Converged Voronoi cells
\begin{scope}[shift={(11.2,0)}]
  % Points
  \foreach \p in {(0.8,3.2),(1.2,2.8),(0.6,2.7),(1.4,3.3),(1.0,3.5)} \node[ptA] at \p {};
  \foreach \p in {(3.2,3.4),(3.8,3.1),(3.3,2.9),(3.7,3.5),(3.5,2.7)} \node[ptB] at \p {};
  \foreach \p in {(2.0,1.1),(2.5,0.8),(1.8,0.7),(2.3,1.3),(2.1,0.5)} \node[ptC] at \p {};

  % Centroids at true centers
  \node[cent, fill=sourceblue!50] (c1) at (1.0, 3.1) {$\mu_1$};
  \node[cent, fill=targetgreen!50] (c2) at (3.5, 3.1) {$\mu_2$};
  \node[cent, fill=quotientorange!50] (c3) at (2.1, 0.9) {$\mu_3$};

  % Converged Voronoi boundaries
  \draw[thick, densely dashed, structurepurple!80!black] (2.25, 4.0) -- (2.25, 2.1) -- (0, 1.8);
  \draw[thick, densely dashed, structurepurple!80!black] (2.25, 2.1) -- (4.2, 1.8);

  \node[lbl, align=center] at (2.1, -0.6) {\textbf{(c) Converged}\\Voronoi boundaries settle\\Cost $J$ minimized};
\end{scope}
\end{tikzpicture}
\caption{Lloyd's $k$-means clustering algorithm in geometric space. \textbf{(a)}
Centroids $\mu_1, \mu_2, \mu_3$ are initialized arbitrarily, partitioning the
points into Voronoi cells. \textbf{(b)} In the update step, each centroid slides
to the mean (center of mass) of its assigned cluster. \textbf{(c)} In the
assignment step, points rejoin the nearest updated centroid; boundaries settle
cleanly between natural clusters.}
\label{ch36:fig-kmeans-loop}
\end{figure}

\subsection{The cost that $k$-means minimizes}

What does ``tidier'' mean precisely? We measure how spread out each cluster is.
Let $\mu_1,\dots,\mu_k$ be the centroids and let $c(i)\in\{1,\dots,k\}$ be the
cluster index assigned to point $x^{(i)}$. The cost is the total squared
distance from points to their own centroid:
\[
  J(c,\mu) \;=\; \sum_{i=1}^{n} \big\lVert x^{(i)} - \mu_{c(i)} \big\rVert^2 .
\]
Here $\lVert v\rVert^2 = v_1^2 + v_2^2 + \cdots + v_d^2$ is the squared length of
the vector $v$. This number is small exactly when every point sits close to its
centroid -- which is what we mean by good clusters.

\begin{focusbox}
$k$-means uses \textbf{squared} distance, not plain distance. That is not an
arbitrary choice. Squared distance is what makes step~2 -- ``put the centroid
at the mean'' -- exactly optimal, as we prove below. Different distance choices
give different algorithms (e.g.\ the \emph{$k$-medoids} or \emph{$k$-medians}
variants).
\end{focusbox}

\subsection{Lloyd's algorithm}

The standard procedure is \textbf{Lloyd's algorithm}. It just alternates the two
easy steps until nothing changes.

\begin{algorithm}[H]
\caption{\textsc{$k$-Means} (Lloyd's algorithm)}
\KwIn{points $x^{(1)},\dots,x^{(n)}\in\R^d$; number of clusters $k$}
\KwOut{centroids $\mu_1,\dots,\mu_k$ and an assignment $c:\{1..n\}\to\{1..k\}$}
initialize $\mu_1,\dots,\mu_k$ (e.g.\ pick $k$ distinct points at random)\;
\Repeat{no assignment changed}{
  \tcp{Assignment step: each point joins its nearest centroid}
  \For{$i \leftarrow 1$ \KwTo $n$}{
    $c(i) \leftarrow \arg\min_{j\in\{1..k\}} \big\lVert x^{(i)} - \mu_j \big\rVert^2$\;
  }
  \tcp{Update step: each centroid moves to the mean of its members}
  \For{$j \leftarrow 1$ \KwTo $k$}{
    $S_j \leftarrow \{\, i : c(i) = j \,\}$\;
    \lIf{$S_j \neq \emptyset$}{$\mu_j \leftarrow \dfrac{1}{|S_j|}\sum_{i\in S_j} x^{(i)}$}
  }
}
\Return $(\mu, c)$\;
\end{algorithm}

\subsection{A fully worked numeric run}

Let us cluster six points on the line ($d=1$, so the math is light) into $k=2$
groups. The points are
\[
  x = 1,\;\; 2,\;\; 4,\;\; 7,\;\; 8,\;\; 9 .
\]
By eye there are two clusters, $\{1,2,4\}$ and $\{7,8,9\}$. Watch the algorithm
discover them.

\textbf{Initialize.} Suppose our random pick lands on the two centroids
$\mu_1 = 2$ and $\mu_2 = 8$ (the lab below traces an unfavorable initial pair
so you can see how initialization affects the outcome).

\textbf{Round 1, assignment.} For each point pick the nearer centroid (distance,
not squared, gives the same winner):

\begin{center}
\begin{tabular}{ccccc}
\toprule
point & $|x-2|$ & $|x-8|$ & nearer & cluster \\
\midrule
$1$ & $1$ & $7$ & $\mu_1$ & $1$ \\
$2$ & $0$ & $6$ & $\mu_1$ & $1$ \\
$4$ & $2$ & $4$ & $\mu_1$ & $1$ \\
$7$ & $5$ & $1$ & $\mu_2$ & $2$ \\
$8$ & $6$ & $0$ & $\mu_2$ & $2$ \\
$9$ & $7$ & $1$ & $\mu_2$ & $2$ \\
\bottomrule
\end{tabular}
\end{center}

So $S_1 = \{1,2,4\}$ and $S_2 = \{7,8,9\}$.

\textbf{Round 1, update.}
\[
  \mu_1 = \frac{1+2+4}{3} = \frac{7}{3} \approx 2.33,
  \qquad
  \mu_2 = \frac{7+8+9}{3} = 8 .
\]

\textbf{Round 2, assignment.} Re-check every point against the new centroids
$2.33$ and $8$. Point $4$ is at distance $1.67$ from $\mu_1$ and $4$ from
$\mu_2$, so it stays in cluster~1. Every point keeps its cluster: \emph{no
assignment changed}. The loop stops.

\textbf{Final cost.}
\[
  J = \underbrace{(1-\tfrac73)^2 + (2-\tfrac73)^2 + (4-\tfrac73)^2}_{\text{cluster 1}}
    + \underbrace{(7-8)^2 + (8-8)^2 + (9-8)^2}_{\text{cluster 2}}
    = \tfrac{14}{3} + 2 \approx 6.67.
\]
The algorithm recovered exactly the clustering you would have drawn by hand,
and it did so in two passes.

\begin{center}
\begin{tikzpicture}[scale=0.95]
  % number line
  \draw[->] (0,0) -- (10.5,0);
  \foreach \x in {1,2,4,7,8,9} \draw (\x,0.08) -- (\x,-0.08) node[below]{\footnotesize \x};
  % cluster 1 points
  \foreach \x in {1,2,4} \fill[sourceblue] (\x,0) circle (2.6pt);
  % cluster 2 points
  \foreach \x in {7,8,9} \fill[targetgreen] (\x,0) circle (2.6pt);
  % centroids
  \draw[quotientorange,thick] (2.333,0.45) node[above]{\footnotesize $\mu_1=\frac73$} -- (2.333,0.05);
  \draw[quotientorange,thick] (8,0.45) node[above]{\footnotesize $\mu_2=8$} -- (8,0.05);
  \node[sourceblue] at (2.4,-0.75) {\footnotesize cluster 1};
  \node[targetgreen] at (8,-0.75) {\footnotesize cluster 2};
\end{tikzpicture}
\end{center}

\subsection{Why does it stop? (Correctness as a decreasing potential)}

How do we know Lloyd's algorithm does not loop forever, ping-ponging between
configurations? Because the cost $J$ acts like a \emph{potential} that can only
go down, and there are only finitely many possible assignments.

\begin{lemma}[Each step does not increase the cost]
\label{ch36:lemma-monotone}
In Lloyd's algorithm, both the assignment step and the update step leave
$J(c,\mu)$ unchanged or make it smaller.
\end{lemma}

\begin{focusbox}
The proof rests on two ``easy optimization'' facts:
\begin{itemize}
  \item \textbf{Assignment step.} With the centroids $\mu$ held \emph{fixed},
        $J = \sum_i \lVert x^{(i)} - \mu_{c(i)}\rVert^2$ is minimized over the
        choice of $c(i)$ by sending each point to its nearest centroid -- the
        terms do not interact, so you minimize each one separately.
  \item \textbf{Update step.} With the assignment $c$ held \emph{fixed}, the
        best single point to represent a set $S$ under \emph{squared} distance
        is its mean. This is the fact below.
\end{itemize}
\end{focusbox}

\begin{lemma}[The mean minimizes squared distance]
\label{ch36:lemma-mean}
For any finite set of points $\{z_1,\dots,z_m\}\subset\R^d$, the function
$g(\mu) = \sum_{t=1}^{m}\lVert z_t - \mu\rVert^2$ is minimized at
$\mu^\star = \frac{1}{m}\sum_t z_t$, the mean.
\end{lemma}

\begin{proof}
Work one coordinate at a time; the coordinates separate because the squared
norm is a sum over coordinates. In one dimension we minimize
$g(\mu)=\sum_t (z_t-\mu)^2$. This is a parabola opening upward in $\mu$. Take
the derivative and set it to zero:
\[
  g'(\mu) = \sum_{t=1}^{m} 2(\mu - z_t) = 2\Big(m\mu - \sum_t z_t\Big) = 0
  \;\;\Longrightarrow\;\;
  \mu^\star = \frac{1}{m}\sum_{t} z_t .
\]
Since $g''(\mu) = 2m > 0$, this critical point is the unique minimum. Doing this
in every coordinate gives the vector mean.
\end{proof}

\begin{proof}[Proof of Lemma~\ref{ch36:lemma-monotone}]
The assignment step minimizes $J$ over $c$ with $\mu$ fixed, so it cannot raise
$J$. The update step replaces each centroid by the mean of its cluster, which by
Lemma~\ref{ch36:lemma-mean} minimizes that cluster's contribution to $J$ with
$c$ fixed, so it cannot raise $J$ either.
\end{proof}

\begin{theorem}[Lloyd's algorithm terminates]
\label{ch36:thm-terminate}
Lloyd's algorithm halts after finitely many rounds.
\end{theorem}

\begin{proof}
There are at most $k^n$ distinct assignments $c$. Each round either changes the
assignment -- in which case, by Lemma~\ref{ch36:lemma-monotone}, the cost
strictly drops (an assignment changes only when a point finds a strictly closer
centroid) -- or it does not, in which case the loop stops. Because the cost
strictly decreases on every productive round and the cost is determined by the
assignment, no assignment can ever recur. With finitely many assignments, the
process must reach a round where nothing changes.
\end{proof}

\begin{trapbox}
Termination is \emph{not} the same as optimality. Lloyd's algorithm finds a
\emph{local} minimum -- a configuration no single re-assignment or re-centering
can improve -- but possibly not the \emph{global} best clustering. A bad random
start can strand it in a poor local minimum. In fact, finding the globally
optimal $k$-means clustering is NP-hard. This is the same greedy/local-search
limitation you saw in earlier chapters.
\end{trapbox}

\begin{labbox}
\textbf{Trace a bad start.} Re-run the six-point example but initialize
$\mu_1 = 1$, $\mu_2 = 2$ (both centroids inside the left clump). Do the
assignment step (point $1$ to $\mu_1$; everything $\ge 2$ to $\mu_2$, since it is
closer to $2$ than to $1$), recompute the means, and keep going. Where does it
settle? Compare its final cost to the $6.67$ we got above. Then try
$\mu_1 = 2,\ \mu_2 = 4$. You will see why practitioners run
$k$-means several times from different starts and keep the lowest-cost result.
This restart trick, when the starts are chosen cleverly, is called
\textbf{$k$-means++}.
\end{labbox}

\begin{complexitybox}
\textbf{Cost of one round of Lloyd's algorithm.} The assignment step compares
each of $n$ points against each of $k$ centroids, each comparison touching $d$
coordinates: $\Th(nkd)$. The update step sums each point into one centroid:
$\Th(nd)$. So a round is $\Th(nkd)$, dominated by assignment.

If the algorithm runs for $T$ rounds, the total is $\Th(Tnkd)$. In the worst
case $T$ can be exponential, but in practice $T$ is small (often tens of
rounds), so $k$-means is treated as a near-linear-time algorithm. Memory is
$\Th(nd + kd)$.
\end{complexitybox}

\begin{aestheticbox}
$k$-means is everywhere because it is cheap and intuitive: compressing an image
to $k$ colors (\emph{color quantization}), grouping customers for marketing,
summarizing documents, initializing more elaborate models, and compressing
vectors for fast nearest-neighbor search. When an engineer says ``let us just
cluster it and see,'' they almost always mean run $k$-means.
\end{aestheticbox}

\section{Gradient descent: walking downhill on purpose}
\label{ch36:gradient-descent}

$k$-means minimized its cost by a clever alternation that happened to have a
closed-form best move at each step. Most learning problems are not so lucky:
there is no formula for the best $\theta$, only a cost surface $J(\theta)$ we
can evaluate and \emph{differentiate}. For those, the standard tool
is \textbf{gradient descent}.

\subsection{The gradient is the uphill arrow}

Recall from calculus that the derivative $J'(\theta)$ of a one-variable function
tells you its slope: if $J'(\theta) > 0$ the function is rising as $\theta$
increases, so to go \emph{down} you should \emph{decrease} $\theta$. To descend,
step \emph{against} the sign of the derivative.

In many variables the derivative becomes the \textbf{gradient}, written
$\nabla J(\theta)$ -- the vector of partial derivatives, one per knob:
\[
  \nabla J(\theta) =
  \left(\frac{\partial J}{\partial \theta_1},\,
        \frac{\partial J}{\partial \theta_2},\,\dots,\,
        \frac{\partial J}{\partial \theta_p}\right).
\]

\begin{intuitionbox}
At any point on the cost surface, the gradient $\nabla J(\theta)$ points in the
direction of \emph{steepest ascent} -- the compass bearing that climbs fastest.
You want to go down, so you step in the \emph{opposite} direction,
$-\nabla J(\theta)$. ``Gradient descent'' literally
means ``repeatedly step a little bit in the most-downhill direction.''
\end{intuitionbox}

\subsection{The update rule and the learning rate}

You take steps of size controlled by a number $\eta > 0$ (Greek ``eta''), the
\textbf{learning rate} or \textbf{step size}:
\[
  \theta \;\leftarrow\; \theta \;-\; \eta\,\nabla J(\theta).
\]
You repeat this until the gradient is nearly zero (the ground is flat) or you run
out of patience.

\begin{algorithm}[H]
\caption{\textsc{Gradient-Descent}}
\KwIn{differentiable cost $J:\R^p\to\R$; learning rate $\eta>0$; iterations $T$}
\KwOut{a parameter vector $\theta$ with (we hope) small $J(\theta)$}
initialize $\theta$ (e.g.\ all zeros or small random values)\;
\For{$t \leftarrow 1$ \KwTo $T$}{
  compute the gradient $g \leftarrow \nabla J(\theta)$\;
  $\theta \leftarrow \theta - \eta\, g$\;
  \lIf{$\lVert g\rVert$ is tiny}{\textbf{break} \tcp*[f]{ground is flat: stop}}
}
\Return $\theta$\;
\end{algorithm}

\begin{trapbox}
The learning rate $\eta$ is the number beginners get wrong most often.
\begin{itemize}
  \item Too \textbf{small}: you inch downhill and it takes forever (and may stall
        before reaching the valley within your iteration budget).
  \item Too \textbf{large}: you overshoot the valley, land on the far slope
        \emph{higher} than you started, overshoot back, and \emph{diverge} --
        the cost explodes to infinity.
\end{itemize}
Choosing $\eta$ means finding the stable middle range between those two regimes. We will see the divergence concretely below.
\end{trapbox}

\subsection{A fully worked one-variable run}

Take the simplest non-trivial cost, a parabola:
\[
  J(\theta) = (\theta - 3)^2, \qquad J'(\theta) = 2(\theta - 3).
\]
The minimum occurs at $\theta = 3$ where $J = 0$. Let us pretend we do not
know that and let gradient descent find it. Start at $\theta_0 = 0$ with
$\eta = 0.1$. The update is
$\theta \leftarrow \theta - 0.1\cdot 2(\theta-3) = \theta - 0.2(\theta-3)
 = 0.8\,\theta + 0.6.$

\begin{center}
\begin{tabular}{cccc}
\toprule
step $t$ & $\theta_t$ & $J'(\theta_t)=2(\theta_t-3)$ & $J(\theta_t)$ \\
\midrule
$0$ & $0.000$ & $-6.000$ & $9.000$ \\
$1$ & $0.600$ & $-4.800$ & $5.760$ \\
$2$ & $1.080$ & $-3.840$ & $3.686$ \\
$3$ & $1.464$ & $-3.072$ & $2.359$ \\
$4$ & $1.771$ & $-2.458$ & $1.510$ \\
$5$ & $2.017$ & $-1.966$ & $0.966$ \\
$\vdots$ & $\vdots$ & $\vdots$ & $\vdots$ \\
$10$ & $2.678$ & $-0.644$ & $0.104$ \\
$20$ & $2.965$ & $-0.069$ & $0.0012$ \\
\bottomrule
\end{tabular}
\end{center}

Every step, the derivative is negative (we are to the left of
the minimum, slope going down-to-the-right), so the update \emph{adds} to
$\theta$, nudging us rightward toward $3$. The cost falls steadily. The steps
shrink as we approach, because the gradient itself shrinks near the bottom --
so gradient descent automatically takes smaller steps near a minimum. After 20
steps we are at $2.965$, very close to the minimum.

\begin{center}
\begin{tikzpicture}[scale=0.95]
  \draw[->] (-0.3,0) -- (6.3,0) node[right]{\footnotesize $\theta$};
  \draw[->] (0,-0.3) -- (0,5.2) node[above]{\footnotesize $J$};
  % parabola J=(theta-3)^2 scaled by 0.5 vertically to fit
  \draw[sourceblue,thick,domain=0.05:5.95,samples=80]
       plot (\x, {0.5*(\x-3)*(\x-3)});
  \node[sourceblue] at (5.2,4.4) {\footnotesize $J(\theta)=(\theta-3)^2$};
  % iterates as descending dots
  \foreach \x in {0,0.6,1.08,1.464,1.771,2.017}{
     \fill[trapred] (\x, {0.5*(\x-3)*(\x-3)}) circle (2.2pt);
  }
  \draw[->,quotientorange,thick] (0,4.5) .. controls (1,2.5) and (2,1.0) .. (2.95,0.05);
  \node[quotientorange] at (1.0,3.3) {\footnotesize steps};
  \fill (3,0) circle (1.5pt) node[below]{\footnotesize $\theta^\star=3$};
\end{tikzpicture}
\end{center}

\begin{labbox}
\textbf{Make it diverge.} Re-run the same parabola from $\theta_0 = 0$ but set
$\eta = 1.1$. Now the update is $\theta \leftarrow \theta - 1.1\cdot 2(\theta-3)
= -1.2\,\theta + 6.6$. Compute $\theta_1, \theta_2, \theta_3$. You will get
$6.600$, then $-1.320$, then $8.184$ -- the iterates fling \emph{outward},
alternating sides of $3$ and getting worse each time. The cost blows up. This is what
``$\eta$ too large'' looks like. Find by experiment the
largest $\eta$ that still converges. (For $J=(\theta-3)^2$ the boundary is
$\eta = 1$; convergence needs $\eta < 1$.)
\end{labbox}

\subsection{Convexity: when downhill is enough}

Gradient descent found the bottom of the parabola because a parabola has exactly
one valley. A cost function with a single bowl shape -- no false valleys, no
plateaus that trap you -- is called \textbf{convex}. On convex costs, ``keep
walking downhill'' is guaranteed to reach the true global minimum.

\begin{formalbox}
A function $J:\R^p\to\R$ is \textbf{convex} if for all $\theta,\theta'$ and all
$\lambda\in[0,1]$,
\[
  J\big(\lambda\theta + (1-\lambda)\theta'\big)
  \;\le\;
  \lambda J(\theta) + (1-\lambda) J(\theta').
\]
Geometrically: the straight line segment joining any two points on the graph
never dips \emph{below} the graph. Such a function has no local minimum that is
not also global. The squared-error costs of $k$-means' update step and of linear
models (next section) are convex in their parameters; this is why they are easy
to optimize.
\end{formalbox}

Figure~\ref{ch36:fig-convex} makes the contrast visible. On the left, a convex
bowl: the chord between any two points stays \emph{above} the curve, and wherever
you start, downhill leads to the single global minimum. On the right, a wavy
non-convex cost: the chord dips below the curve in places, and which valley you
fall into now depends on where you begin.

\begin{figure}[H]
\centering
\begin{tikzpicture}[scale=0.9]
  % ---- convex bowl ----
  \begin{scope}
    \draw[->] (-0.3,0) -- (4.3,0) node[right]{\footnotesize $\theta$};
    \draw[->] (0,-0.3) -- (0,3.2) node[above]{\footnotesize $J$};
    \draw[sourceblue,thick,domain=0.2:3.8,samples=60]
         plot (\x,{0.35*(\x-2)*(\x-2)+0.3});
    % chord above the curve
    \fill (0.6,{0.35*(0.6-2)*(0.6-2)+0.3}) circle (1.5pt);
    \fill (3.4,{0.35*(3.4-2)*(3.4-2)+0.3}) circle (1.5pt);
    \draw[trapred,dashed] (0.6,{0.35*(0.6-2)*(0.6-2)+0.3}) -- (3.4,{0.35*(3.4-2)*(3.4-2)+0.3});
    \node[trapred] at (2.5,2.6) {\footnotesize chord};
    \fill[targetgreen] (2,0.3) circle (2pt) node[below]{\footnotesize global min};
    \node[sourceblue] at (2,-0.7) {\footnotesize convex: one valley};
  \end{scope}
  % ---- non-convex ----
  \begin{scope}[xshift=6.5cm]
    \draw[->] (-0.3,0) -- (4.3,0) node[right]{\footnotesize $\theta$};
    \draw[->] (0,-0.3) -- (0,3.2) node[above]{\footnotesize $J$};
    \draw[structurepurple,thick,domain=0.2:3.8,samples=80]
         plot (\x,{0.33*(\x-2)*(\x-2)*(\x-2)*(\x-2) - 1.4*(\x-2)*(\x-2) + 2.2});
    \fill[targetgreen] (0.54,{0.33*(0.54-2)*(0.54-2)*(0.54-2)*(0.54-2) - 1.4*(0.54-2)*(0.54-2) + 2.2}) circle (2pt);
    \fill[targetgreen] (3.46,{0.33*(3.46-2)*(3.46-2)*(3.46-2)*(3.46-2) - 1.4*(3.46-2)*(3.46-2) + 2.2}) circle (2pt);
    \node[structurepurple] at (2,-0.7) {\footnotesize non-convex: many valleys};
  \end{scope}
\end{tikzpicture}
\caption{Convex versus non-convex cost surfaces. On a convex cost (left) every
chord lies above the curve and downhill always reaches the single global minimum.
On a non-convex cost (right) the two marked valleys are different local minima;
the starting point decides which one gradient descent finds.}
\label{ch36:fig-convex}
\end{figure}

\begin{walkbox}
Consider this question while you stretch your legs. Convexity guarantees a
single global valley, so gradient descent on a convex cost cannot get stuck in a
bad local minimum. Yet the modern deep-learning models you hear about are
\emph{not} convex -- their loss surfaces are complicated, full of valleys. They still
work surprisingly well. So: what does ``learning succeeded'' really mean if we
never reach the global minimum? Is a ``good enough'' valley the real goal? Consider
the possibility that for these problems, \emph{finding any deep-enough
valley} is the primary objective, and the global minimum may not even be desirable
(it might just memorize the data).
\end{walkbox}

\subsection{Batch vs.\ stochastic: the great speedup}

When the cost is an average over $n$ data points,
$J(\theta) = \frac1n\sum_i \ell_i(\theta)$, the gradient is also an average:
$\nabla J(\theta) = \frac1n\sum_i \nabla \ell_i(\theta)$. Computing it exactly
means touching all $n$ examples for \emph{every single step}. With millions of
examples that is expensive.

\begin{intuitionbox}
You do not need the exact downhill direction -- you only need a direction that is
\emph{downhill on average}. So estimate the gradient from a small random sample
of the data (even a single example!) and step using that noisy estimate. Each
step is now cheap, so you take many more of them. The noise can even help you
escape shallow local minima. This is \textbf{stochastic gradient descent (SGD)},
the method behind essentially all modern large-scale learning.
\end{intuitionbox}

Figure~\ref{ch36:fig-sgd} shows the trade in one picture, looking down on the cost
surface from above (the closed curves are contours of equal cost, like a
topographic map). Batch descent takes a few smooth, expensive steps straight for
the minimum. SGD jitters along a noisier path -- each step points only roughly
downhill -- but every step is so cheap that it covers far more ground per unit of
work.

\begin{figure}[H]
\centering
\begin{tikzpicture}[scale=0.85]
  % concentric contour ellipses (a bowl seen from above)
  \foreach \r in {0.6,1.2,1.8,2.4}{
    \draw[gray!55] (0,0) ellipse ({1.3*\r} and {0.85*\r});
  }
  \fill[targetgreen] (0,0) circle (2pt) node[below=2pt,font=\footnotesize]{min};
  % batch path: smooth, few steps
  \coordinate (b0) at (-3.3,1.9);
  \draw[->,sourceblue,very thick] (b0) -- (-1.9,1.0) -- (-0.9,0.45) -- (-0.25,0.12) -- (0,0);
  \fill[sourceblue] (b0) circle (2pt);
  \node[sourceblue,font=\footnotesize] at (-3.3,2.3) {batch start};
  % SGD path: jittery, many short steps
  \coordinate (s0) at (3.3,-1.9);
  \draw[->,trapred,thick]
    (s0) -- (2.6,-1.0) -- (2.4,-1.4) -- (1.7,-0.6) -- (1.6,-0.95)
        -- (1.0,-0.2) -- (0.8,-0.55) -- (0.35,0.05) -- (0.2,-0.25) -- (0,0);
  \fill[trapred] (s0) circle (2pt);
  \node[trapred,font=\footnotesize] at (3.3,-2.3) {SGD start};
\end{tikzpicture}
\caption{Two paths down the same bowl, drawn as a contour map. Batch gradient
descent (blue) takes a few exact steps; stochastic gradient descent (red) takes
many cheap, noisy steps. Per step SGD wanders, but per unit of computation it
usually reaches the valley faster on large datasets.}
\label{ch36:fig-sgd}
\end{figure}

\begin{complexitybox}
Let one gradient term $\nabla\ell_i$ cost $\Th(d)$ to evaluate.
\begin{itemize}
  \item \textbf{Batch} gradient descent: $\Th(nd)$ per step, $T$ steps,
        $\Th(Tnd)$ total. Each step is an \emph{exact} downhill move.
  \item \textbf{Stochastic} gradient descent with a minibatch of size $b$:
        $\Th(bd)$ per step. With $b=1$ that is $\Th(d)$ per step -- independent
        of dataset size $n$. The trade is noisier steps, so you usually need
        more of them, but each is so cheap that SGD wins decisively at scale.
\end{itemize}
\end{complexitybox}

\section{The perceptron: learning a straight-line boundary}
\label{ch36:perceptron}

Now we put the pieces together on a \emph{supervised} problem and meet one of
the earliest learning algorithms. The \textbf{perceptron} (Rosenblatt, 1958) is
a learning rule for telling two classes apart with a flat boundary. It is the
ancestor of every neural network, and it is simple enough to trace by hand.

\subsection{The setup: a linear classifier}

Each example is a feature vector $x\in\R^d$ with a label $y\in\{-1,+1\}$ (think
``cat vs.\ not-cat,'' ``spam vs.\ not-spam''). The model is a weight vector
$w\in\R^d$ and a bias $b\in\R$. It predicts
\[
  \hat{y}(x) = \operatorname{sign}\big(w\cdot x + b\big),
  \qquad w\cdot x = \sum_{j=1}^{d} w_j x_j .
\]
The set of points where $w\cdot x + b = 0$ is a flat boundary -- a line in $2$D,
a plane in $3$D, a \textbf{hyperplane} in general. Points on one side get $+1$,
the other side $-1$. Training means \emph{finding} a $(w,b)$ whose boundary puts
the two classes on opposite sides.

\begin{focusbox}
A convenient trick removes the bias from the bookkeeping: glue a constant $1$
onto every input, $\tilde{x} = (x_1,\dots,x_d,1)$, and fold $b$ into the weights
as the last coordinate, $\tilde{w} = (w_1,\dots,w_d,b)$. Then
$w\cdot x + b = \tilde{w}\cdot\tilde{x}$, and we only ever update one vector.
We use this trick in the trace below.
\end{focusbox}

\subsection{The learning rule}

The perceptron's update is very simple. Loop through the examples.
When the current weights classify an example correctly, do nothing. When they get
one \emph{wrong}, nudge the weights toward fixing that example:
\[
  \text{if } y^{(i)}\,(w\cdot x^{(i)}) \le 0 \;\;(\text{a mistake}),\qquad
  w \leftarrow w + y^{(i)}\, x^{(i)} .
\]

\begin{algorithm}[H]
\caption{\textsc{Perceptron} (with the bias folded into $w$)}
\KwIn{examples $(x^{(1)},y^{(1)}),\dots,(x^{(n)},y^{(n)})$ with $y^{(i)}\in\{-1,+1\}$; max passes $T$}
\KwOut{weight vector $w$ separating the classes (if one exists)}
$w \leftarrow \mathbf{0}$\;
\For{$\text{pass} \leftarrow 1$ \KwTo $T$}{
  $\textit{mistakes} \leftarrow 0$\;
  \For{$i \leftarrow 1$ \KwTo $n$}{
    \If{$y^{(i)}\,(w\cdot x^{(i)}) \le 0$}{
      $w \leftarrow w + y^{(i)}\, x^{(i)}$\;
      $\textit{mistakes} \leftarrow \textit{mistakes} + 1$\;
    }
  }
  \lIf{$\textit{mistakes} = 0$}{\textbf{break} \tcp*[f]{a full clean pass: done}}
}
\Return $w$\;
\end{algorithm}

\begin{intuitionbox}
Why does adding $y^{(i)} x^{(i)}$ fix a mistake? Suppose the true label is $+1$
but $w\cdot x^{(i)}$ came out $\le 0$ (wrong). We add $x^{(i)}$ to $w$. The new
score on that same point becomes
$(w + x^{(i)})\cdot x^{(i)} = w\cdot x^{(i)} + \lVert x^{(i)}\rVert^2$, which is
\emph{larger} than before -- pushed in the $+$ direction, exactly where we wanted
it. For a $-1$ example we subtract $x^{(i)}$ and the score drops. Every update
moves the boundary so the offending point is ``more correct.''
\end{intuitionbox}

\subsection{A fully worked trace}

Let us separate two classes in the plane. Using the bias trick, each input is
$(x_1,x_2,1)$ and $w=(w_1,w_2,w_3)$. Four examples:

\begin{center}
\begin{tabular}{ccc}
\toprule
$i$ & $x^{(i)}=(x_1,x_2,1)$ & $y^{(i)}$ \\
\midrule
$1$ & $(2,\,1,\,1)$ & $+1$ \\
$2$ & $(0,\,-1,\,1)$ & $-1$ \\
$3$ & $(-1,\,0,\,1)$ & $-1$ \\
$4$ & $(1,\,1,\,1)$ & $+1$ \\
\bottomrule
\end{tabular}
\end{center}

Start $w=(0,0,0)$. With $w=\mathbf 0$ the score is $0$, and $y\cdot 0 = 0 \le 0$,
so the very first point counts as a mistake.

\textbf{Example 1} $(2,1,1)$, $y=+1$: score $w\cdot x = 0 \le 0$ -- mistake.
Update $w \leftarrow (0,0,0) + (+1)(2,1,1) = (2,1,1)$.

\textbf{Example 2} $(0,-1,1)$, $y=-1$: score
$(2,1,1)\cdot(0,-1,1) = 0 - 1 + 1 = 0$, and $y\cdot 0 = 0 \le 0$ -- mistake.
Update $w \leftarrow (2,1,1) + (-1)(0,-1,1) = (2,2,0)$.

\textbf{Example 3} $(-1,0,1)$, $y=-1$: score
$(2,2,0)\cdot(-1,0,1) = -2 + 0 + 0 = -2$, and $y\cdot(-2) = (-1)(-2) = 2 > 0$ --
correct! No update.

\textbf{Example 4} $(1,1,1)$, $y=+1$: score
$(2,2,0)\cdot(1,1,1) = 2 + 2 + 0 = 4$, and $y\cdot 4 = 4 > 0$ -- correct! No
update.

That ends pass~1 with $w=(2,2,0)$ and two mistakes. Pass~2: re-check all four.

\begin{center}
\begin{tabular}{ccccc}
\toprule
$i$ & $x^{(i)}$ & $y^{(i)}$ & score $w\cdot x$ & $y\cdot\text{score}$ \\
\midrule
$1$ & $(2,1,1)$ & $+1$ & $4+2+0=6$ & $6>0$ correct \\
$2$ & $(0,-1,1)$ & $-1$ & $0-2+0=-2$ & $2>0$ correct \\
$3$ & $(-1,0,1)$ & $-1$ & $-2+0+0=-2$ & $2>0$ correct \\
$4$ & $(1,1,1)$ & $+1$ & $2+2+0=4$ & $4>0$ correct \\
\bottomrule
\end{tabular}
\end{center}

Zero mistakes on a full pass -- we stop. The final classifier is
$w=(2,2,0)$, i.e.\ predict $\operatorname{sign}(2x_1 + 2x_2)$, whose boundary is
the line $x_1 + x_2 = 0$. Check it against the data: it really does put the two
$+$ points on one side and the two $-$ points on the other.

Figure~\ref{ch36:fig-perceptron} illustrates both the geometric vector update
and the final converged separating boundary.

\begin{figure}[H]
\centering
\begin{tikzpicture}[
  scale=0.92,
  >=Stealth,
  lbl/.style={font=\sffamily\footnotesize}
]
% Panel (a): Geometric update mechanism
\begin{scope}[shift={(0,0)}]
  \draw[->, gray!60] (-2.2,0) -- (3.5,0) node[right, lbl, black]{$x_1$};
  \draw[->, gray!60] (0,-2.2) -- (0,2.6) node[above, lbl, black]{$x_2$};

  % Initial weight vector w_old
  \coordinate (Wold) at (0.4, 1.8);
  \draw[->, very thick, sourceblue] (0,0) -- (Wold);
  \node[above left=2pt, lbl, sourceblue] at (Wold) {$\mathbf{w}_{\text{old}}$};

  % Old boundary: perpendicular to w_old
  \draw[dashed, sourceblue!70, thick] (-2.2, 0.49) -- (2.2, -0.49);
  \node[above right, lbl, sourceblue!70, font=\sffamily\scriptsize] at (-2.2, 0.52) {$\mathbf{w}\cdot x = 0$};

  % Misclassified point x^(i) in negative halfspace
  \coordinate (X) at (1.8, 0.4);
  \fill[targetgreen!80!black] (X) circle (2.8pt);
  \node[above right=2pt, lbl, targetgreen!80!black] at (X) {$x^{(i)}$ ($y=+1$)};

  % Vector addition: w_new = w_old + x^(i)
  \coordinate (Wnew) at ($(Wold) + (X)$);
  \draw[->, thick, dashed, targetgreen!80!black] (Wold) -- (Wnew)
    node[midway, above=4pt, lbl, targetgreen!80!black]{$+x^{(i)}$};
  \draw[->, very thick, trapred!80!black] (0,0) -- (Wnew)
    node[above right=2pt, lbl, trapred!80!black]{$\mathbf{w}_{\text{new}}$};

  % New tilted boundary
  \draw[very thick, trapred!70, domain=-1.8:2.2] plot (\x, {-(\x*2.2)/2.2});

  \node[lbl, align=center] at (0.5, -3.1) {\textbf{(a) Weight Vector Update}\\[3pt]Adding $y^{(i)}x^{(i)}$ tilts $\mathbf{w}$\\[3pt]Hyperplane rotates toward $x^{(i)}$};
\end{scope}

% Panel (b): Converged classifier on the 4 trace points
\begin{scope}[shift={(7.5,0)}]
  \draw[->, gray!60] (-2.2,0) -- (3.2,0) node[right, lbl, black]{$x_1$};
  \draw[->, gray!60] (0,-2.2) -- (0,2.6) node[above, lbl, black]{$x_2$};

  % Final boundary line x1 + x2 = 0
  \draw[very thick, structurepurple, domain=-2.0:2.0] plot (\x, {-\x});
  \node[structurepurple, lbl, above] at (-1.85, 2.02) {$x_1 + x_2 = 0$};

  % Normal vector w = (2,2) drawn from the boundary in upper-left quadrant
  \draw[->, very thick, sourceblue] (-0.7, 0.7) -- ++(0.9, 0.9)
    node[above right=2pt, lbl, sourceblue]{$\mathbf{w} = (2,2)$};
  % right-angle indicator at (-0.7, 0.7)
  \draw[gray!70] (-0.7, 0.7) ++(0.12, 0.12) -- ++(0.12, -0.12) -- ++(-0.12, -0.12);

  % Half-plane shading indications
  \node[targetgreen!80!black, lbl] at (1.8, 2.2) {\scriptsize $\mathbf{w}\cdot x > 0$ ($+$ side)};
  \node[trapred!80!black, lbl] at (-1.4, -1.8) {\scriptsize $\mathbf{w}\cdot x < 0$ ($-$ side)};

  % Positive points
  \fill[targetgreen!80!black] (2, 1) circle (2.6pt) node[right=2pt, lbl, targetgreen!80!black]{$(2,1)$ [$+$]};
  \fill[targetgreen!80!black] (1, 1) circle (2.6pt) node[below right=2pt, lbl, targetgreen!80!black]{$(1,1)$ [$+$]};

  % Negative points
  \fill[trapred!80!black] (0, -1) circle (2.6pt) node[left=3pt, lbl, trapred!80!black]{$(0,-1)$ [$-$]};
  \fill[trapred!80!black] (-1, 0) circle (2.6pt) node[below left=2pt, lbl, trapred!80!black]{$(-1,0)$ [$-$]};

  \node[lbl, align=center] at (0.5, -3.1) {\textbf{(b) Converged Decision Boundary}\\[3pt]Zero errors on all $4$ points\\[3pt]Normal $\mathbf{w}$ points into positive class};
\end{scope}
\end{tikzpicture}
\caption{The perceptron decision boundary and update geometry. \textbf{(a)} When
a positive example $x^{(i)}$ is misclassified, the update $\mathbf{w} \leftarrow \mathbf{w} + x^{(i)}$
tilts the weight vector toward $x^{(i)}$, rotating the orthogonal hyperplane
decision boundary $\mathbf{w}\cdot x = 0$ until $x^{(i)}$ lands on the positive side.
\textbf{(b)} The converged classifier on the four trace points: boundary line
$x_1 + x_2 = 0$ perfectly separates the positive points from the negative
points, with normal vector $\mathbf{w}=(2,2)$ pointing into the positive half-space.}
\label{ch36:fig-perceptron}
\end{figure}

Notice the weight vector $w=(2,2)$ (dropping the bias coordinate) points
\emph{perpendicular} to the boundary, toward the positive side. That is no
accident: $w$ is always the normal direction of the hyperplane it defines.

\subsection{Why it works: the perceptron convergence theorem}

The key fact, and the reason the perceptron caused such excitement, is
that if the data \emph{can} be separated by a hyperplane at all, this simple
update rule is \emph{guaranteed} to find one -- in a bounded number of mistakes
that does not depend on how many examples there are.

\begin{formalbox}
The data is \textbf{linearly separable with margin} $\gamma > 0$ if there is a
unit weight vector $w^\star$ ($\lVert w^\star\rVert = 1$) such that every example
satisfies $y^{(i)}(w^\star\cdot x^{(i)}) \ge \gamma$. Let
$R = \max_i \lVert x^{(i)}\rVert$ be the radius of the data.
\end{formalbox}

\begin{theorem}[Perceptron convergence; Novikoff, 1962]
\label{ch36:thm-perceptron}
If the data is linearly separable with margin $\gamma$ and radius $R$, the
perceptron algorithm makes at most $\big(R/\gamma\big)^2$ mistakes before it
classifies every example correctly (and then halts).
\end{theorem}

\begin{proof}
We track two quantities about $w$ as mistakes accumulate; let $w_k$ be the
weights after the $k$-th mistake, starting from $w_0=\mathbf 0$.

\emph{Lower bound (the weights line up with $w^\star$).} On a mistake using
example $(x,y)$ we set $w_k = w_{k-1} + y\,x$. Dot with $w^\star$:
\[
  w^\star\cdot w_k = w^\star\cdot w_{k-1} + y\,(w^\star\cdot x)
  \;\ge\; w^\star\cdot w_{k-1} + \gamma,
\]
using the margin condition. After $k$ mistakes, $w^\star\cdot w_k \ge k\gamma$.

\emph{Upper bound (the weights cannot grow too fast).} Compute the squared
length after the update:
\[
  \lVert w_k\rVert^2 = \lVert w_{k-1}\rVert^2 + 2y\,(w_{k-1}\cdot x) + \lVert x\rVert^2 .
\]
Because we only update on a \emph{mistake}, we have $y\,(w_{k-1}\cdot x)\le 0$,
and $\lVert x\rVert^2 \le R^2$, so
$\lVert w_k\rVert^2 \le \lVert w_{k-1}\rVert^2 + R^2$. After $k$ mistakes,
$\lVert w_k\rVert^2 \le k R^2$.

\emph{Squeeze.} By the Cauchy--Schwarz inequality (and $\lVert w^\star\rVert=1$),
$w^\star\cdot w_k \le \lVert w_k\rVert$. Combine the two bounds:
\[
  k\gamma \;\le\; w^\star\cdot w_k \;\le\; \lVert w_k\rVert \;\le\; \sqrt{k}\,R .
\]
Thus $k\gamma \le \sqrt{k}\,R$, i.e.\ $\sqrt{k}\le R/\gamma$, giving
$k \le (R/\gamma)^2$. The number of mistakes is bounded, so after finitely many
the algorithm makes a clean pass and stops.
\end{proof}

\begin{focusbox}
The bound $(R/\gamma)^2$ says nothing about $n$, the number of
examples, and nothing about $d$, the dimension. It depends only on the
\emph{geometry}: how big the data is ($R$) versus how comfortably it can be
split ($\gamma$). Wide margin $\Rightarrow$ few mistakes $\Rightarrow$ fast
learning. The margin later became the central quantity of
\emph{support vector machines}, the dominant classifiers of the 1990s and 2000s.
\end{focusbox}

\begin{trapbox}
The convergence theorem assumes the data \emph{is} separable. If it is not -- if
no straight boundary can split the classes -- the perceptron never makes a clean
pass and loops until it hits $T$, with the weights never settling. Historically,
Minsky and Papert (1969) pointed out that even the tiny XOR
function is not linearly separable, so a single perceptron \emph{cannot} learn
it. That critique froze the field for years. The fix -- stacking perceptrons
into layers and training them with gradient descent -- is the neural network,
and it is why Section~\ref{ch36:gradient-descent} matters so much.
\end{trapbox}

\subsection{The perceptron as gradient descent}

The two halves of this chapter are connected. The perceptron update
\emph{is} a stochastic gradient step on a particular cost. Define the per-example
\textbf{perceptron loss}
\[
  \ell(w; x,y) = \max\big(0,\; -\,y\,(w\cdot x)\big),
\]
which is $0$ when the example is classified correctly (with positive score) and
grows linearly with how wrong you are otherwise. Its gradient with respect to
$w$ is $0$ on a correct example and $-\,y\,x$ on a mistake. So a stochastic
gradient \emph{descent} step with learning rate $1$,
$w \leftarrow w - \eta\,\nabla\ell = w + y\,x$ on mistakes, is exactly
Rosenblatt's rule. Thus the 1958 perceptron update is stochastic subgradient
descent applied to hinge-at-zero loss.

\begin{aestheticbox}
Everything in this chapter is used in systems you rely on daily. A spam
filter is a linear classifier trained with SGD. The recommendation feed that
clusters users and items leans on $k$-means-style grouping. A large language
model is a very large stack of differentiable layers whose billions of parameters
are tuned by exactly the gradient-descent loop in
Section~\ref{ch36:gradient-descent} -- just with a more elaborate $J$ and a lot
more hardware. The basic loop is the same in all of them: define a cost, take
the gradient, step downhill, repeat.
\end{aestheticbox}

\section{Summary}
\label{ch36:connections}

The methods in this chapter connect to much of the earlier material:

\begin{itemize}
  \item \textbf{Greedy and local search.} Lloyd's algorithm is alternating
        minimization: each step takes the locally best assignment or center,
        eventually settling into a local minimum.
  \item \textbf{Loop invariants and potentials.} We proved $k$-means halts using
        a strictly decreasing potential $J$, the same potential argument used
        for amortized analysis. We proved perceptron convergence with a two-sided
        squeeze on weight vector growth.
  \item \textbf{Recurrences and cost analysis.} We derived per-round costs
        $\Th(nkd)$ and $\Th(Tnd)$ from loop structure and matrix operations.
  \item \textbf{Linear algebra and geometry.} Dot products, hyperplanes, margins,
        and projections are the native geometric language of data and parameters.
  \item \textbf{Tradeoffs and limits.} Real-world learning balances optimization
        against generalization. Non-convex loss surfaces can trap descent in poor
        local valleys; unregularized models risk memorizing noise rather than
        patterns; and statistical models lack the hard worst-case correctness
        guarantees of classical algorithms.
\end{itemize}

In theory, training a machine-learning model is pure continuous optimization over
linear algebra. In practice, computing gradients across billions of parameters
strains physical hardware to its limits. Arithmetic throughput is rarely the true
bottleneck; the limiting factor is moving bits across the memory bus, cache
hierarchies, and DRAM banks. To understand why algorithms run fast or slow on
real hardware, we need to look at the machine itself. In
Chapter~37, we step from algorithmic models into systems
architecture and memory physics, before exploring how optimizing compilers bridge
source code to machine instructions (Chapter~38) and how
quantum mechanics changes what can be computed efficiently (Chapter~39).

\begin{recoverybox}
The one idea to remember is this loop:
\emph{(1) write a number that measures wrongness; (2) find the direction that
makes it smaller; (3) take a small step; (4) repeat.} Clustering, classification,
regression, and modern neural networks are all variations on
those four lines. Algorithms that learn are still algorithms; what drives them is a cost
function.
\end{recoverybox}

\section*{Exercises}
\label{ch36:exercises}
\addcontentsline{toc}{section}{Exercises}

\begin{exercise}[Warmup: assign and update]
\label{ch36:ex-kmeans-1d}
Run one full round of Lloyd's algorithm by hand on the one-dimensional points
$0, 1, 10, 11, 12$ with $k=2$ and initial centroids $\mu_1 = 0$, $\mu_2 = 1$.
Show the assignment, then the updated centroids. Run a second round. Does the
assignment change? Report the final clustering and its cost $J$.
\end{exercise}

\begin{exercise}[The mean, again]
\label{ch36:ex-mean}
Prove that for points $z_1,\dots,z_m$ on the real line, the value $\mu$ that
minimizes $\sum_t |z_t - \mu|$ (note: \emph{absolute} value, not squared) is a
\emph{median}, not the mean. Explain in one sentence why $k$-means therefore uses
the mean while $k$-medians uses the median.
\end{exercise}

\begin{exercise}[Gradient descent on a parabola]
\label{ch36:ex-gd-parabola}
For $J(\theta) = (\theta-5)^2$ with start $\theta_0 = 0$ and learning rate
$\eta = 0.25$, write the closed-form update $\theta_{t+1} = a\theta_t + c$ and
compute $\theta_1,\theta_2,\theta_3$. Then find, by solving the recurrence, a
formula for $\theta_t$ and confirm $\theta_t \to 5$.
\end{exercise}

\begin{exercise}[The edge of stability]
\label{ch36:ex-gd-diverge}
For $J(\theta) = (\theta - c)^2$, show that the gradient-descent update can be
written $\theta_{t+1} - c = (1 - 2\eta)(\theta_t - c)$. Conclude that the
iterates converge to $c$ if and only if $|1 - 2\eta| < 1$, i.e.\ $0 < \eta < 1$,
and that $\eta = 0.5$ jumps to the exact minimum in one step. Explain the
behavior at $\eta = 1$.
\end{exercise}

\begin{exercise}[Trace the perceptron]
\label{ch36:ex-perceptron-trace}
Using the bias trick (append a $1$ to each input), run the perceptron from
$w=\mathbf 0$ on: $(1,0;+1)$, $(0,1;+1)$, $(-1,-1;-1)$ (here the semicolon
separates features from label). Show every score and every update. How many
mistakes occur before a clean pass? Draw the final boundary.
\end{exercise}

\begin{exercise}[Why $y\,x$?]
\label{ch36:ex-perceptron-fix}
Suppose example $(x,y)$ with $y=+1$ is currently misclassified, so
$w\cdot x \le 0$. After the update $w' = w + x$, show the new score increased by
exactly $\lVert x\rVert^2$. Then show that for a $y=-1$ mistake the update
$w' = w - x$ \emph{decreases} the score by $\lVert x\rVert^2$. Why does this not
guarantee the point is now correctly classified after a single update?
\end{exercise}

\begin{exercise}[Margin and mistakes]
\label{ch36:ex-margin}
Two datasets are each linearly separable. Dataset A has radius $R=2$ and margin
$\gamma = 1$; dataset B has radius $R = 10$ and margin $\gamma = 0.5$. Use
Theorem~\ref{ch36:thm-perceptron} to bound the number of perceptron mistakes on
each. Which dataset is ``harder'' for the perceptron, and by what factor?
\end{exercise}

\begin{exercise}[XOR is not linearly separable]
\label{ch36:ex-xor}
The XOR data is $(0,0;-1)$, $(1,1;-1)$, $(1,0;+1)$, $(0,1;+1)$. Suppose a
separating $w=(w_1,w_2)$ and bias $b$ existed. Write the four inequalities
$y^{(i)}(w\cdot x^{(i)} + b) > 0$ and derive a contradiction. Conclude that no
single perceptron can learn XOR.
\end{exercise}

\begin{exercise}[Batch vs.\ stochastic cost]
\label{ch36:ex-sgd}
A dataset has $n = 10^6$ examples in $d = 100$ dimensions. Estimate the number of
multiply-add operations for one \emph{batch} gradient step versus one
\emph{stochastic} step with minibatch size $b = 32$. If a fixed compute budget
allows $10^9$ multiply-adds, roughly how many steps of each can you afford?
\end{exercise}

\begin{exercise}[Challenge: a non-convex trap]
\label{ch36:ex-nonconvex}
Consider $J(\theta) = \theta^4 - 4\theta^2 + \theta$. Find $J'(\theta)$ and argue
(by examining sign changes) that $J$ has two local minima. Pick two starting
points for gradient descent that you predict will converge to \emph{different}
minima, and explain your reasoning geometrically. This is a simple example of
why neural-network training depends on initialization.
\end{exercise}

\begin{exercise}[Challenge: $k$-means is a coordinate descent]
\label{ch36:ex-coord}
Explain why Lloyd's algorithm can be viewed as \emph{alternating minimization}:
it minimizes $J(c,\mu)$ over $c$ with $\mu$ fixed, then over $\mu$ with $c$
fixed, forever alternating. Using this view, give a one-paragraph argument for
why the cost sequence $J_0 \ge J_1 \ge J_2 \ge \cdots$ is monotone and bounded
below by $0$, hence convergent. Why does convergence of the \emph{cost} not by
itself prove the \emph{assignment} stops changing?
\end{exercise}
